En un ensayo simulado, un depósito de capacidad \(80\,\mathrm L\) contiene inicialmente \(50\,\mathrm L\) de agua limpia. Entran \(3\,\mathrm{L/min}\) de agua limpia y sale \(1\,\mathrm{L/min}\) de mezcla. El contenido se mantiene uniforme y el soluto no reacciona ni se pierde por otros mecanismos. Un dosificador añade soluto sin aportar un volumen apreciable; puede suministrar entre \(0\) y \(2\,\mathrm{g/min}\).

Se desea que la concentración siga la ley \(c(t)=0.02t\,\mathrm{g/L}\), donde \(t\) es el valor numérico del tiempo transcurrido en minutos. Elija las cantidades necesarias para describir el proceso y formule sus balances. Determine cómo debe variar el suministro del dosificador y durante cuánto tiempo puede mantenerse la concentración solicitada, respetando tanto la capacidad del depósito como la del dosificador. Indique cuál de las dos limita primero la operación.
